\documentclass[conference]{IEEEtran}
\IEEEoverridecommandlockouts
\usepackage{cite}
\usepackage{amsmath,amssymb,amsfonts}
\usepackage{algorithmic}
\usepackage{graphicx}
\usepackage{textcomp}
\usepackage{xcolor}
\def\BibTeX{{\rm B\kern-.05em{\sc i\kern-.025em b}\kern-.08em
    T\kern-.1667em\lower.7ex\hbox{E}\kern-.125emX}}
\begin{document}

\title{Knowledge Graph-Augmented Retrieval-Augmented Generation for ISRO Domain Question Answering on Resource-Constrained Hardware}

\author{\IEEEauthorblockN{Nandana Narayan Das}
\IEEEauthorblockA{\textit{Alliance School of Advanced Computing} \\
\textit{Alliance University}\\
Bengaluru, India \\
nandanadas007@gmail.com}
\and
\IEEEauthorblockN{Gowri Kannan}
\IEEEauthorblockA{\textit{Dept. of Computer Science and Engineering} \\
\textit{Viswajyothi College of Engineering and Technology}\\
Muvattupuzha, India \\
gowrikannan479@gmail.com}
}

\maketitle

\begin{abstract}
Retrieval-Augmented Generation (RAG) is a practical approach for domain-specific question answering, but government corpora such as ISRO documentation pose additional challenges because the data are heterogeneous, relational, and often require temporal reasoning. We present KG-RAG, a locally run pipeline that builds a domain knowledge graph from unstructured ISRO documents and combines it with FAISS retrieval and a local Ollama model. We evaluate the system on a deterministic benchmark of 200 ISRO questions with a canonical 20-dev / 180-test split. On the final 180-question test set, BM25 + LLM attains ROUGE-L 0.2915 and reference-token coverage 0.4340, Vanilla RAG attains 0.2780 and 0.3989, and KG-RAG attains 0.2736 and 0.3921. These results are lexical metrics and do not establish factual correctness. They suggest that KG augmentation can help in targeted relational and temporal settings, but improvements are task- and tier-specific under local inference constraints. The system is designed for resource-constrained deployment and does not require cloud-hosted model access.
\end{abstract}

\begin{IEEEkeywords}
Knowledge Graph, Retrieval-Augmented Generation, ISRO, Named Entity Recognition, Mistral-7B, FAISS, Domain-Specific QA, Resource-Constrained Inference
\end{IEEEkeywords}

\section{Introduction}

The Indian Space Research Organisation (ISRO) has been producing mission reports, press releases, payload specifications, and scientific literature for over six decades. This wealth of information is spread across thousands of unstructured web pages, making it difficult to retrieve precise factual answers — especially those that require understanding the relationships between missions, satellites, launch vehicles, and scientists. Standard keyword search fails at these relational queries, and general-purpose language models frequently hallucinate ISRO-specific facts simply because the domain is underrepresented in their training data.

Retrieval-Augmented Generation \cite{lewis2020rag} addresses hallucination by fetching relevant documents at query time and grounding the model's response in retrieved text. However, RAG treats all retrieved content as a flat collection of passages, with no awareness of the structured relationships those documents contain. Knowledge Graph-augmented RAG systems \cite{he2024gretriever, edge2024graphrag, guo2024lightrag, ma2024tog2} take this further by incorporating entity-relation triples alongside passage retrieval, which helps considerably on multi-hop questions. The problem is that every published KG-RAG system has been built around pre-existing structured knowledge bases and tested only on general-domain benchmarks — and almost all require cloud-hosted LLMs to function.

We address two gaps that remain unresolved in the literature: no KG-RAG system exists for any government space agency corpus, and no published system can build its knowledge graph from raw unstructured documents without cloud infrastructure. Our contributions are:

\begin{enumerate}
    \item An engineering integration of local NLP tools (spaCy, NetworkX, FAISS, Ollama) into a domain-specific KG-RAG pipeline for ISRO documentation, showing that effective domain QA does not require cloud infrastructure or pre-built knowledge bases.
    \item \textbf{ISRO-QA}: a benchmark of 200 manually verified question-answer pairs for evaluating space agency domain QA systems.
    \item A fully reproducible, zero-cost deployment on commodity hardware using only open-source tools.
    \item Ablation experiments that isolate and measure how much the knowledge graph contributes over dense retrieval alone.
\end{enumerate}

\section{Related Work}
\label{sec:related}

\subsection{Foundational RAG Architectures}

Lewis et al. \cite{lewis2020rag} established the RAG paradigm by pairing a Dense Passage Retrieval module with a BART generator over a FAISS-indexed Wikipedia corpus, showing that fetching relevant documents at inference time allows a model to answer factual questions more reliably than relying on memorized training knowledge. The retrieved passages are treated as an unordered bag of text with no notion of which entities are related or how. Sarthi et al. \cite{sarthi2024raptor} later improved retrieval granularity through RAPTOR, which recursively summarizes document clusters into a hierarchical tree, achieving around 20\% improvement on long-document QA. RAPTOR is effective but computationally expensive and still does not capture relational structure between named entities.

\subsection{Knowledge Graph Integration with Language Models}

Pan et al. \cite{pan2024unifying} surveyed KG-LLM integration across three paradigms — using KGs to enrich LLMs, using LLMs to build KGs, and tightly coupled systems — and specifically noted that no synergized KG+LLM system had been validated on government or scientific corpora where the KG must be built without manual annotation. Liang et al. \cite{liang2024kag} built KAG for legal and financial QA by combining manually annotated domain KGs with semantic reasoning over Neo4j. KAG works well in those domains, but its reliance on human schema design makes it impractical to deploy in a new technical domain like ISRO without substantial expert involvement.

\subsection{Graph-Based RAG Systems}

Several systems have explored using knowledge graphs as the retrieval backbone instead of flat vector indices. GraphRAG \cite{edge2024graphrag} builds a KG using an LLM and applies Leiden community detection to generate multi-level summaries. LightRAG \cite{guo2024lightrag} achieves comparable results more cheaply by combining entity-level and concept-level retrieval. HippoRAG \cite{gutierrez2024hipporag} uses Personalized PageRank over an open KG to simulate associative memory retrieval, outperforming dense baselines on MuSiQue and HotpotQA. All three outperform vanilla RAG on standard multi-hop benchmarks, but all depend on GPT-4 class models for KG construction, have never been tested on government scientific corpora, and assume well-structured large document collections.

\subsection{Multi-Hop Reasoning with Knowledge Graphs}

Think-on-Graph 2.0 \cite{ma2024tog2} improves multi-hop accuracy by alternating Freebase KG lookups with unstructured passage retrieval, gaining 8\% on WebQSP and CWQ. GNN-RAG \cite{mavromatis2024gnnrag} trains a graph neural network on KG subgraphs to identify answer paths, achieving 85.7 hits@1 on WebQSP. GRAG \cite{hu2024grag} reduces irrelevant graph context through soft subgraph pruning and graph-aware prompting. Wang et al. \cite{wang2024kgprompting} dynamically construct a KG from input documents using GPT-4, improving multi-document QA F1 by 4.2\% over RAG. All of these approaches assume either Freebase, Wikidata, or a cloud LLM — none provide a way to build the required KG from raw public documents alone.

\subsection{Advanced KG-RAG Frameworks}

G-Retriever \cite{he2024gretriever} formulates subgraph extraction as a Prize-Collecting Steiner Tree problem and generates answers from the extracted subgraph. KG$^2$RAG \cite{zhang2025kg2rag} combines entity linking with KG neighbourhood traversal and vector retrieval, outperforming LightRAG and GraphRAG on multi-hop benchmarks. Tomczak et al. \cite{tomczak2025kgextended} add question decomposition and chain-of-thought reasoning on top of Wikidata-backed retrieval. Dong et al. \cite{dong2025kaqg} ground question generation in KG-extracted facts for difficulty-controlled educational QA. All of these systems either require pre-existing structured KBs or GPT-4 for graph construction.

\subsection{Domain-Specific RAG and Evaluation}

Graber et al. \cite{graber2025docgraphrag} showed that adapting GraphRAG to manufacturing — building KGs from technical manuals — consistently outperforms naive RAG. Rakin et al. \cite{rakin2024domainrag} found that fine-tuning RAG on domain data cuts hallucination rates from 29\% to 6\%. KNOWNET \cite{deng2024knownet} integrates GPT-3.5 with UMLS and MeSH health KGs to make medical retrieval safer. These papers confirm that domain adaptation pays off, but all require supervised fine-tuning datasets or pre-curated KBs that do not exist for ISRO. RAGAS \cite{es2023ragas} provides automated reference-free scoring via Faithfulness, Answer Relevancy, and Context Relevance metrics, validated on Wikipedia-domain annotations. Peng et al. \cite{peng2024graphragsurvey} surveyed five Graph RAG paradigms and explicitly flagged the absence of scientific or government domain benchmarks as an open problem, which directly motivated us to build ISRO-QA.

\section{System Architecture and Methodology}
\label{sec:method}

Figure~\ref{fig:architecture} shows the KG-RAG pipeline, organized around four stages: document ingestion, offline indexing, query-time hybrid retrieval, and answer generation. Two indexing branches run in parallel — one builds a knowledge graph from entity-relation triples extracted by spaCy, and the other encodes document chunks into a FAISS vector index using MiniLM-L6-v2. At query time, both branches contribute context that is merged and passed to a locally hosted Mistral-7B model. Figure~\ref{fig:architecture} makes clear that the entire system operates without any external API calls — every component from NLP processing to final answer generation runs on the local machine.

\begin{figure}[h]
\centering
\includegraphics[width=\columnwidth]{kg_rag_system_architecture.png}
\caption{KG-RAG system architecture: document ingestion, KG construction (2,193 nodes, 4,655 edges), FAISS indexing (4,557 chunks, 384-dim), hybrid retrieval, and local LLM generation via Ollama.}
\label{fig:architecture}
\end{figure}

\subsection{Document Collection and Preprocessing}

We crawled \texttt{isro.gov.in} using the Firecrawl API, which handles JavaScript rendering and PDF extraction, returning clean markdown per page. The crawler followed links up to three hops from seed pages covering missions, annual reports, press releases, launch vehicles, and personnel profiles. Each document was split into overlapping 512-token chunks with a 128-token stride to prevent sentences from being cut at boundaries. Navigation menus, footers, and boilerplate text were stripped using a custom cleaning script. The final corpus contains 339 documents and 4,557 chunks averaging 9,874 characters each, with no empty or malformed chunks.

\subsection{Knowledge Graph Construction}

We used spaCy's \texttt{en\_core\_web\_lg} pipeline to extract entities and build the knowledge graph. NER identifies five categories of domain-relevant entities: missions (e.g., Chandrayaan-2, Mangalyaan), satellites and payloads (e.g., RISAT-2BR1, SAR payload), launch vehicles (e.g., PSLV-C37, GSLV Mk III), scientists and personnel (e.g., S. Somanath, K. Sivan), and temporal expressions such as launch dates. For relation extraction, we used spaCy's dependency parser to find the syntactic head connecting two co-occurring named entities in the same sentence. When two entities are linked by a dependency path of at most two hops — chosen to capture direct subject-verb-object relations without introducing noise from longer chains — we extract a triple of the form \textit{(subject, relation, object)}, where the relation label comes from the arc type and governing verb. For instance, \textit{``Chandrayaan-2 was launched by ISRO in July 2019''} produces \texttt{(Chandrayaan-2, launched\_by, ISRO)} and \texttt{(Chandrayaan-2, launch\_date, July 2019)}. All triples are stored in a NetworkX directed multigraph, giving a graph of 2,193 nodes and 4,655 typed edges. We note that spaCy's NER was not trained on ISRO-specific terminology and misses a number of domain-specific entity variants; a custom entity ruler is a priority for future work.

\subsection{Dense Vector Indexing}

Every preprocessed chunk is encoded using \texttt{all-MiniLM-L6-v2}, producing 384-dimensional embeddings. We ran the encoder on CPU rather than GPU to keep the full VRAM budget available for Mistral at inference time. The embeddings are loaded into a FAISS Flat L2 index, which performs exact nearest-neighbour search. This approach is appropriate for a corpus of 4,557 chunks; corpora beyond approximately 10,000 chunks would require approximate indexing strategies such as FAISS IVF or HNSW to maintain acceptable query latency. Both the FAISS index and the NetworkX graph are saved to disk after indexing and reloaded at query time.

\subsection{Hybrid Retrieval and Answer Generation}

When a query arrives, we first apply a keyword pre-filter — chunks containing none of the query keywords are excluded before dense ranking. The threshold of 20 matching chunks before falling back to full FAISS search was set empirically on 20 held-out development questions. We then run spaCy NER on the query to extract entity names and retrieve all one-hop neighbours from the knowledge graph along with their edge labels, serializing these as natural language triples to form the graph context. Simultaneously, we encode the query and retrieve the top-$k$ nearest chunks from FAISS — with $k$=3, chosen to balance context length against retrieval diversity — as the passage context. The two streams are joined (graph triples first, then passages) and fed into Mistral-7B-Instruct (Q4\_K\_M) via Ollama, which fits in approximately 4.1~GB of VRAM. The system prompt instructs the model to answer strictly from context and say ``I don't know'' when evidence is insufficient. Zero-shot inference is used throughout with no fine-tuning.

\section{Dataset and Experimental Setup}
\label{sec:experiments}

\subsection{ISRO-QA Benchmark}

We manually built ISRO-QA, a benchmark of 200 question-answer pairs, each verified against at least one document in the corpus. Reference answers were either copied verbatim from the source or paraphrased closely to preserve factual fidelity. All 200 pairs were curated by the primary author and cross-verified against source documents by the second author; inter-annotator agreement on answer correctness was not formally measured and constitutes a limitation of the current benchmark. Questions fall into three tiers: 100 factoid questions with single entity answers such as launch dates and payload names, 60 multi-hop relational questions that connect facts from two or more documents, and 40 timeline reasoning questions requiring duration computation or event ordering. The full benchmark is released alongside the code.

\subsection{Baselines and Evaluation Metrics}

We compare KG-RAG against two baselines under identical hardware and corpus conditions. BM25+LLM retrieves the top passages using BM25 sparse retrieval and passes them to Mistral-7B for generation, representing traditional keyword-based retrieval. Vanilla RAG swaps BM25 for FAISS dense retrieval but removes the knowledge graph entirely, isolating the KG contribution. Systems are evaluated on ROUGE-L (longest common subsequence F1 against the reference), reference-token coverage, Exact Match, and IDK Rate (fraction of abstentions). We acknowledge that these are lexical metrics that reward surface-form overlap rather than factual correctness; human evaluation is therefore prepared separately and has not been completed. The canonical 180-question test set is frozen; paired confidence intervals and Wilcoxon tests are reported in the additional analysis, without rerunning generation. The system uses spaCy 3.7, NetworkX 3.3, FAISS-CPU 1.8, sentence-transformers, and Ollama.

\section{Results and Discussion}
\label{sec:results}

\subsection{Main Results}

Table~\ref{tab:results} reports the frozen 180-question test set, not the 20-question development set. BM25+LLM is strongest overall on the two principal lexical metrics (ROUGE-L 0.2915; coverage 0.4340). KG-RAG (0.2736; 0.3921) does not outperform BM25 and is slightly below Vanilla RAG (0.2780; 0.3989). The paired analysis finds no significant KG-RAG advantage over Vanilla RAG for ROUGE-L, coverage, exact match, or IDK (all two-sided Wilcoxon $p \geq 0.356$ where defined). Against BM25, KG-RAG has significantly lower coverage ($p=0.011$) and a higher IDK rate ($p<0.001$). These are lexical findings and do not establish factual correctness.

\begin{table}[h]
\centering
\caption{Canonical performance on the frozen ISRO-QA test set (180 questions).}
\label{tab:results}
\begin{tabular}{lcccc}
\hline
\textbf{System} & \textbf{ROUGE-L} & \textbf{Coverage} & \textbf{EM} & \textbf{IDK\%} \\
\hline
BM25 + LLM             & 0.2915 & 0.4340 & 0.0278 & 1.67\% \\
Vanilla RAG            & 0.2780 & 0.3989 & 0.0333 & 12.22\% \\
\textbf{KG-RAG (ours)} & 0.2736 & 0.3921 & 0.0333 & 11.67\% \\
\hline
\end{tabular}
\end{table}

Figure~\ref{fig:results} visualizes the canonical point estimates. The figure should be read as a comparison of lexical overlap and abstention behavior, not as evidence of factual superiority.

\begin{figure}[h]
\centering
\includegraphics[width=\columnwidth]{results_comparison.png}
\caption{ROUGE-L and reference-token coverage on the frozen 180-question test set. BM25+LLM is the strongest lexical performer; KG-RAG is not universally superior.}
\label{fig:results}
\end{figure}

\subsection{Per-Tier Analysis}

Table~\ref{tab:tier} breaks down abstention rates by tier on the test split. KG-RAG has its highest IDK rate on Tier 1 factoid questions (15.56\%) and its lowest on Tier 3 timeline questions (5.56\%). Relative to Vanilla RAG, the largest directional KG improvement is on Tier 3: ROUGE-L is higher by 0.0224 and IDK is lower by 8.33 percentage points, but neither difference is statistically significant in the paired tier analysis ($p=0.485$ and $p=0.083$, respectively). The complete Tier 2 and Tier 3 comparison is shown in Fig.~\ref{fig:tiercomparison}; it should be interpreted as exploratory evidence rather than a claim of significance.

\begin{table}[h]
\centering
\caption{IDK rate (\%) per difficulty tier.}
\label{tab:tier}
\begin{tabular}{lccc}
\hline
\textbf{System} & \textbf{Tier 1 (90)} & \textbf{Tier 2 (54)} & \textbf{Tier 3 (36)} \\
\hline
BM25 + LLM      & 3.33\%  & 0.00\% & 0.00\% \\
Vanilla RAG     & 12.22\%  & 11.11\% & 13.89\% \\
\textbf{KG-RAG} & 15.56\% & 9.26\% & 5.56\% \\
\hline
\end{tabular}
\end{table}

\begin{figure}[h]
\centering
\includegraphics[width=\columnwidth]{tier_comparison.png}
\caption{Tier-wise ROUGE-L and reference-token coverage on the complete canonical test tiers.}
\label{fig:tiercomparison}
\end{figure}

\subsection{Paired Statistical Analysis}

The statistical analysis treats the question as the unit of analysis and uses the same stored answer for every paired comparison. For KG-RAG minus Vanilla RAG, the mean paired differences are $-0.0044$ for ROUGE-L (95\% bootstrap CI $[-0.0277,0.0185]$), $-0.0068$ for coverage (CI $[-0.0334,0.0202]$), and $-0.0056$ for IDK rate (CI $[-0.0444,0.0333]$); the corresponding Wilcoxon $p$-values are 0.403, 0.357, and 0.782. Thus, the frozen benchmark does not support a significant overall KG advantage over Vanilla RAG. Against BM25, KG-RAG's coverage difference is $-0.0420$ (CI $[-0.0724,-0.0121]$, $p=0.011$), while its ROUGE-L difference is $-0.0179$ (CI $[-0.0429,0.0086]$, $p=0.183$). Exact-match testing has too few non-zero paired differences for a meaningful Wilcoxon test in the KG-vs-BM25 comparison and is reported as such in the machine-readable artifact.
\subsection{Targeted Aditya-L1 Evaluation}

To test whether the graph provides a more useful signal when questions explicitly depend on mission entities and relationships, we conducted a separate researcher-constructed 36-question Aditya-L1 evaluation using newly added official ISSDC mission documentation \cite{issdcaditya}. This set is not part of the canonical 180-question benchmark and is therefore reported as a targeted stress test rather than independent external validation. KG-RAG achieved ROUGE-L 0.3859 and coverage 0.5473, compared with 0.2924/0.4897 for Vanilla RAG and 0.2918/0.4767 for BM25+LLM. All systems had zero exact matches.

For the two primary ROUGE-L comparisons, paired Wilcoxon tests remained significant after Holm correction ($p_{adj}=0.0098$ versus Vanilla RAG and $p_{adj}=0.0151$ versus BM25+LLM). The paired effect sizes were $d_z=0.444$ and $d_z=0.324$, respectively. The KG-RAG minus Vanilla RAG bootstrap 95\% CI for the mean difference was $[0.0288,0.1695]$; the KG-RAG minus BM25 CI was $[-0.0057,0.1834]$, so the latter interval does not independently establish a positive mean effect. The measured IDK rate was also lower for KG-RAG (8.33\%) than for Vanilla RAG (25.00\%) and BM25+LLM (22.22\%). These results support a task-conditional benefit of explicit graph augmentation for this targeted entity-relational domain, while remaining subject to benchmark-construction and lexical-metric limitations.

\subsection{Ablation Study}

Table~\ref{tab:ablation} presents the separate existing 50-question sample comparing KG-only retrieval, FAISS-only retrieval, and full KG-RAG. KG-only is weakest (ROUGE-L 0.0927; coverage 0.1078; IDK 76\%), while Full KG-RAG exceeds FAISS-only on ROUGE-L and coverage in this subset. This result is limited to the ablation subset and is not evidence that KG-RAG wins the canonical benchmark. Expanded filtering and hop variants remain pending because no compatible raw answer traces are stored.

\begin{table}[h]
\centering
\caption{Ablation study on 50-question sample.}
\label{tab:ablation}
\begin{tabular}{lccc}
\hline
\textbf{Configuration} & \textbf{ROUGE-L} & \textbf{Coverage} & \textbf{IDK\%} \\
\hline
KG-only              & 0.0927 & 0.1078 & 76.0\% \\
FAISS-only           & 0.2830 & 0.4077 & 34.0\% \\
\textbf{Full KG-RAG} & \textbf{0.2952} & \textbf{0.4564} & \textbf{10.0\%} \\
\hline
\end{tabular}
\end{table}

Figure~\ref{fig:ablation} visualizes the separate subset result; it must not be conflated with the canonical 180-question benchmark.

\begin{figure}[h]
\centering
\includegraphics[width=\columnwidth]{ablation_results.png}
\caption{Ablation results on the separate 50-question sample. The Coverage panel (right) shows the measured +4.87 percentage-point difference between Full KG-RAG and FAISS-only.}
\label{fig:ablation}
\end{figure}

\subsection{Discussion}

Three conclusions stand out. First, BM25+LLM is the strongest overall lexical performer, and KG-RAG does not outperform it on the frozen test set. Second, KG-RAG is directionally better than Vanilla RAG on Tier 3 ROUGE-L and IDK, but the paired tier tests do not establish statistical significance; the evidence supports a targeted-benefit hypothesis, not universal superiority. Third, the separate 50-question ablation indicates that the full hybrid configuration is more useful than KG-only or FAISS-only in that subset, while the expanded component and hop analyses remain pending.

The targeted Aditya-L1 evaluation provides the clearest positive result in this study: KG-RAG substantially outperforms both baselines on the selected mission-focused questions. Because the set was constructed by the researchers after corpus expansion, it should be treated as targeted stress-test evidence rather than an unbiased external validation set. The significant paired ROUGE-L tests therefore support a conditional claim about graph augmentation on entity-relational Aditya-L1 questions, not universal KG-RAG superiority. The additional evaluation also makes the remaining evidence gaps explicit. Retrieval Recall@k/MRR cannot be computed without manually verified chunk relevance judgments. Human factual evaluation, KG-triple verification, a verified unanswerable set, and manually reviewed paraphrases are prepared but not completed. Resource profiling measured only BM25 retrieval on three queries; dense retrieval attempted unavailable Hugging Face network access, and generation, peak RAM/VRAM, and tokens/sec were not measured. The FAISS Flat L2 index does not scale past roughly 10,000 chunks with acceptable latency; larger deployments would need approximate indexing. Future work should complete the annotation packages and run the planned additional-subset generation experiments before making factual or efficiency claims.

\section{Conclusion}
\label{sec:conclusion}

We built KG-RAG, a system that answers factual questions about ISRO by combining an automatically constructed domain-specific knowledge graph with FAISS dense retrieval, running locally without cloud-hosted model access. On the frozen 180-question test set, BM25+LLM is the strongest lexical system (ROUGE-L 0.2915; coverage 0.4340), while KG-RAG obtains 0.2736 and 0.3921. Paired tests do not show a significant overall advantage for KG-RAG over Vanilla RAG, although Tier 3 shows a promising but non-significant directional pattern. A separate 36-question Aditya-L1 targeted evaluation provides evidence for a task-conditional KG benefit: KG-RAG reaches ROUGE-L 0.3859 versus 0.2924 for Vanilla RAG and 0.2918 for BM25+LLM, with Holm-adjusted paired Wilcoxon $p$-values of 0.0098 and 0.0151, respectively. The existing 50-question ablation supports retaining the hybrid design as an experimental configuration, not as proof of main-benchmark superiority. The released scripts, raw stored outputs, explicit pending statuses, and annotation templates provide a reproducible basis for completing factual, retrieval, and resource evaluations without overstating current evidence.

\begin{thebibliography}{20}

\bibitem{lewis2020rag}
P. S. H. Lewis, E. Perez, A. Piktus, F. Petroni, V. Karpukhin, N. Goyal, H. K\"{u}ttler, M. Lewis, W. Yih, T. Rockt\"{a}schel, S. Riedel, and D. Kiela, ``Retrieval-Augmented Generation for Knowledge-Intensive NLP Tasks,'' in \textit{Advances in Neural Information Processing Systems (NeurIPS)}, 2020.

\bibitem{he2024gretriever}
X. He, Y. Tian, Y. Sun, N. V. Chawla, T. Laurent, Y. LeCun, X. Bresson, and B. Hooi, ``G-Retriever: Retrieval-Augmented Generation for Textual Graph Understanding and Question Answering,'' in \textit{Advances in Neural Information Processing Systems (NeurIPS)}, 2024.

\bibitem{edge2024graphrag}
D. Edge, H. Trinh, N. Cheng, J. Bradley, A. Chao, A. Mody, S. Truitt, and J. Larson, ``From Local to Global: A Graph RAG Approach to Query-Focused Summarization,'' \textit{arXiv preprint arXiv:2404.16130}, 2024.

\bibitem{guo2024lightrag}
Z. Guo, L. Xia, Y. Yu, T. Ao, and C. Huang, ``LightRAG: Simple and Fast Retrieval-Augmented Generation,'' \textit{arXiv preprint arXiv:2410.05779}, 2024.

\bibitem{ma2024tog2}
S. Ma, X. Xu, X. Tian, T. Yang, Y. Fang, and C. Xu, ``Think-on-Graph 2.0: Deep and Faithful Large Language Model Reasoning with Knowledge-Guided Retrieval,'' in \textit{Proceedings of the 2025 Annual Conference of the North American Chapter of the Association for Computational Linguistics (NAACL)}, 2025.

\bibitem{es2023ragas}
S. Es, J. James, L. Espinosa-Anke, and S. Schockaert, ``RAGAS: Automated Evaluation of Retrieval Augmented Generation,'' in \textit{Proceedings of the 18th Conference of the European Chapter of the Association for Computational Linguistics (EACL)}, 2024.

\bibitem{sarthi2024raptor}
P. Sarthi, S. Abdullah, A. Tuli, S. Khanna, A. Goldie, and C. D. Manning, ``RAPTOR: Recursive Abstractive Processing for Tree-Organized Retrieval,'' in \textit{Proceedings of the 12th International Conference on Learning Representations (ICLR)}, 2024.

\bibitem{pan2024unifying}
S. Pan, L. Luo, Y. Wang, C. Chen, J. Wang, and X. Wu, ``Unifying Large Language Models and Knowledge Graphs: A Roadmap,'' \textit{IEEE Transactions on Knowledge and Data Engineering}, vol. 36, no. 7, pp. 3580--3599, 2024.

\bibitem{liang2024kag}
L. Liang, M. Sun, G. Gui, W. Zhu, J. Tang, Z. Zhang, Z. Shao, Z. Xu, X. Tian, Y. Shi, X. Yao, J. Yang, L. Cheng, X. Zhao, Z. Fang, X. Tang, Q. Liu, and W. Liu, ``KAG: Boosting LLMs in Professional Domains via Knowledge Augmented Generation,'' \textit{arXiv preprint arXiv:2409.13731}, 2024.

\bibitem{gutierrez2024hipporag}
B. F. Guti\'{e}rrez, Y. Shu, Y. Gu, M. Yasunaga, and Y. Su, ``HippoRAG: Neurobiologically Inspired Long-Term Memory for Large Language Models,'' in \textit{Advances in Neural Information Processing Systems (NeurIPS)}, 2024.

\bibitem{mavromatis2024gnnrag}
C. Mavromatis and G. Karypis, ``GNN-RAG: Graph Neural Retrieval for Large Language Model Reasoning,'' \textit{arXiv preprint arXiv:2405.20139}, 2024.

\bibitem{hu2024grag}
Y. Hu, Z. Lei, Z. Zhang, B. Pan, C. Ling, and L. Zhao, ``GRAG: Graph Retrieval-Augmented Generation,'' \textit{arXiv preprint arXiv:2405.16506}, 2024.

\bibitem{wang2024kgprompting}
Y. Wang, N. Lipka, R. A. Rossi, A. Siu, R. Zhang, and T. Derr, ``Knowledge Graph Prompting for Multi-Document Question Answering,'' in \textit{Proceedings of the 38th AAAI Conference on Artificial Intelligence (AAAI)}, 2024.

\bibitem{zhang2025kg2rag}
Q. Zhang, B. Xu, Y. Luo, X. Zheng, D. Zhan, H. Peng, and P. S. Yu, ``KG\textsuperscript{2}RAG: Knowledge Graph-Guided Retrieval Augmented Generation,'' in \textit{Proceedings of the 2025 Annual Conference of the North American Chapter of the Association for Computational Linguistics (NAACL)}, 2025.

\bibitem{tomczak2025kgextended}
J. Tomczak, T. Zaranis, and A. Papadopoulos, ``Knowledge Graph-Extended RAG for Question Answering,'' \textit{arXiv preprint arXiv:2504.08893}, 2025.

\bibitem{dong2025kaqg}
J. Dong, Z. Jiang, Y. Li, and X. Jiang, ``KAQG: A Knowledge-Graph-Enhanced RAG for Difficulty-Controlled Question Generation,'' \textit{arXiv preprint arXiv:2505.07618}, 2025.

\bibitem{graber2025docgraphrag}
M. Graber, J. Schlatt, and A. Karaoglu, ``Document GraphRAG: Knowledge Graph Enhanced RAG for Document Question Answering in the Manufacturing Domain,'' \textit{Electronics}, vol. 14, no. 11, p. 2102, 2025.

\bibitem{rakin2024domainrag}
S. Rakin, M. Islam, and A. Rahman, ``Leveraging Domain Adaptation of Retrieval Augmented Generation Models for Question Answering and Reducing Hallucination,'' \textit{arXiv preprint arXiv:2410.17783}, 2024.

\bibitem{deng2024knownet}
C. Deng, Z. Tian, Z. Liu, P. Chen, S. Wen, and T. Wang, ``KNOWNET: Guided Health Information Seeking from LLMs via Knowledge Graph Integration,'' in \textit{Proceedings of the IEEE Visualization and Visual Analytics (VIS)}, 2024.

\bibitem{issdcaditya}
Indian Space Research Organisation, ``Aditya-L1 Mission and Payload Documentation,'' Indian Space Science Data Centre (ISSDC), 2023--2024.

\bibitem{peng2024graphragsurvey}
B. Peng, Y. Zhu, Y. Liu, X. Bo, H. Shi, C. Hong, Y. Tang, and S. He, ``Graph Retrieval-Augmented Generation: A Survey,'' \textit{arXiv preprint arXiv:2408.08921}, 2024.

\end{thebibliography}
\end{document}
